\section{No-fault, cost, invalid-output, and operational constraints}
\label{sec:constraints}

For a policy $d$ and stratum $j\in\{0,1\}$, define
\begin{equation}
 V_j(d)=\E[Y_i^{d(Z_i,M_i)}\mid J_i=j],\qquad
 G_j(d)=V_j(d)-V_j(g_0).
 \label{eq:stratum-values}
\end{equation}
No-fault noninferiority is
\begin{equation}
 V_0(d)\ge V_0(g_0)-\eta.
 \label{eq:noninferiority}
\end{equation}
This is average no-fault performance degradation, not an individually harmed
fraction.  The condition $J=0$ means that the registered fault was not injected;
it does not guarantee baseline success.  New invalidity, unnecessary
abstention, latency, unwanted state change, and escalation are reported as
separate no-fault endpoints.

Cost remains a vector containing input and output tokens, API expense, CPU and
GPU time, wall time, verifier and tool calls, storage and network use, human
escalation or adjudication time, and artifact rewrites.  Of these, the eight
constrained components $k$ are input and output tokens, API expense, CPU and
GPU seconds, wall time, external calls (with verifier, tool, and controller
calls collapsed into a single component), and human minutes; storage, network
use, and artifact rewrites are reported audit fields rather than constrained
components.  This is what fixes $K=8$ and hence the certification family and
sample ceiling of \cref{eq:combined-family,eq:sample-ceiling}.  For each
component $k$, write the deployment-weighted expected cost using the same
prevalence mixture as \eqref{eq:prevalence-value},
\begin{equation}
 C_k(d;p)=p\,\E\big[C^{d(Z_i,M_i)}_{ik}\mid J_i=1\big]
 +(1-p)\,\E\big[C^{d(Z_i,M_i)}_{ik}\mid J_i=0\big],
 \label{eq:cost-mixture}
\end{equation}
and impose the population constraint
\begin{equation}
 C_k(d;p)\le B_k=\beta_k\,C_k(g_0;p),
 \label{eq:cost-constraint}
\end{equation}
so that each budget is stated as a multiple $\beta_k$ of what the no-intervention
regime spends on the same component at the same prevalence, rather than as an
absolute figure that would not survive a change of workload.  Because costs are
nonnegative, a component on which $g_0$ spends nothing constrains every
candidate to zero in that component as well; no additive floor is inferred.
A hard per-regime
execution ceiling applies as an operational feasibility rule only, and is not
interchangeable with expected population budgets; no scalar cost is defined until
empirical cost distributions are observed and exchange rates approved.  Hard
ceilings are safety rails, not treatment definitions: their target binding rate
is exactly zero, they are not part of the estimand, and the binding rate is
reported per regime in E0, the pilot, and the confirmatory run whatever the
outcome.  A ceiling that binds even occasionally would bind disproportionately on
recovery, the intrinsically most expensive regime, mechanically biasing the
verification-versus-recovery comparison this study exists to measure; any nonzero
binding rate is therefore a reported protocol deviation, never a post-hoc ceiling
increase.

Policy-level accounting charges the decision procedure once: monitor and
availability evaluation, gate arithmetic, and their calls, compute, and latency
are allocated into the existing eight constrained components rather than added
as a ninth.  The formal comparator is a monitored policy that always continues,
so the decision procedure's overhead appears on both sides of
\eqref{eq:cost-constraint}; the literal unmonitored variant is retained as a
descriptive operational-cost sensitivity.  Each $\beta_k$ is justified under
that convention before the contract is frozen.

Which components can bind is a property of the bound substrate, and the registry
is deliberately larger than any one substrate needs.  Under the \emph{signed}
tabular cost contract three constrained components are structurally zero for
every regime: API expense, because there is no provider; human minutes, because
there is no human in the loop; and accelerator seconds, because that substrate
is CPU-only.  Their constraints are vacuously satisfied under the zero-baseline
rule and can never bind.  Under the local-GPU successor the third of these
becomes live --- serving the model locally makes GPU seconds non-zero on every
branch --- and its ceiling is then derived from the substrate's own soak and
enforced like any other live component.  The successor cost contract that makes
that substitution is prepared and unsigned; until it is signed, the signed
contract's three structural zeros are what the guard enforces.  A hosted endpoint would invert the first case in the same way,
with API expense falling hardest on the most expensive regime.  Either
substitution changes the character of the certification family rather than a
value in it, which is precisely why these components are retained rather than
dropped: the registry must not depend on a
substrate that happens not to incur API, GPU, or human cost.  A paper about when
to interrupt an agent prices human time at zero only because every regime is
machine-executed and ``interruption'' means controller intervention, not a human
call-out.  The cost of retaining $K=8$ is at most twelve roots out of the sample
ceiling, because $N$ enters that bound only through a logarithm.

Invalid-output rates are constrained analogously,
\begin{equation}
 R(d;p)=p\E[I_i^{d(Z_i,M_i)}\mid J_i=1]
 +(1-p)\E[I_i^{d(Z_i,M_i)}\mid J_i=0]
 \le \xi.
 \label{eq:invalid-constraint}
\end{equation}
Here $\xi$ is a single tolerance on the prevalence-weighted invalid-output rate.
Stratum-specific tolerances $\xi_0,\xi_1$ may additionally be registered, in
which case they are separate constraints and enter the multiplicity count of
\cref{prop:certification} as such; the mixture bound does not imply them and is
not implied by them.

The values $\eta=0.05$ and $\xi=0.02$ remain planning values and cannot be
selected for power.  The approved cost contract instead fixes the monitored
comparator and $\beta_k=1.25$.  Signed structural ceilings are substrate-bound:
a replacement must reapprove its own before E0, and the generic $1.5\times$
value is inactive.  \Cref{tab:ceiling} shows what each margin costs in roots --- the
dependence on $\eta$ is quadratic, so setting it is a sample-size decision as
much as a substantive one, and precisely why it may not be set by looking at
the sample.

Constraint \eqref{eq:cost-constraint} is evaluated on cache-neutralised resource
components; a locally served substrate reports the corresponding local meter, and
a hosted one reports the provider meter beside it.  Where inference is hosted, a
shared provider cache can make billed cost depend on branch order
(\cref{sec:cache}), which is why the constraint is not written on billed cost.
The illustrative billing calculation is a design calculation, not a measurement
from this study.

One consequence must be stated rather than left implicit.  A per-branch ceiling
on a constrained component is enforced by the budget guard, which terminates the
branch in the budget-exhaustion category, and that category scores zero
completion.  A binding per-branch expense ceiling would therefore depress the
measured success of the most expensive regime through a billing mechanism, in
exactly the way a binding wall-clock rail would through a timing one.  Such
ceilings are safety rails with a target binding rate of zero, derived from each
regime's own structural maximum; the binding rate is reported per regime, and a
non-zero rate is a protocol deviation rather than a data point --- never resolved
by raising the ceiling after the fact.
